\subsection{The last ordinary act profile at eighteen arrows}

\begin{proposition}[Proved: six outside elements over three rows]
\label{CP2:A-three-row-six-outside-acts}
If a finite monoid has minimal band $3\times2$ and
at most six outside elements, every finite right-act
cone has zero rational degree-two products.
\end{proposition}
\begin{proof}
The case of at most five outside elements is
Proposition~\ref{CP2:A-three-row-five-outside-acts}.
Its preceding affine argument applies at any outside
count: a possible nonzero product forces row values
$(0,0,1)$ and the four distinct outside profiles
\[
 \begin{aligned}
 p_0&=((0,0,2),(0,1)),& p_1&=((1,1,2),(0,1)),\\
 r_0&=((2,2,0),(1,0)),& r_1&=((2,2,1),(1,0)).
 \end{aligned}
\]
Together with the identity these occupy five of the
six outside labels. The sixth has one of exactly $24$
balanced profiles: four row maps over identity columns,
two over swapped columns, and nine over each of the
two constant column maps. For a constant column map,
all three row images must lie in one of the two value
levels, giving $2^3+1=9$ choices.

We give a complete small literal calculation for this
last outside label. Fix identity0, the four forced labels
1 through4, the additional label5, and six row-major
band labels6 through11. All products involving the
identity or band are forced. For every other ordered
pair, its allowable outputs are exactly all twelve
labels with the composite row and column profile.
The new producer \path{converse/scope18/forced_three_row12.py}
tries every such output, propagating only required
associativity equalities. Every monoid under consideration
is obtained after the indicated relabelling: no ordinary
isomorphism or opposition identification removes a table.
There are exactly $27$ associative tables, reached in
$44$ search nodes across the $24$ extra profiles.

The independent program
\path{converse/scope18/check_forced_three_row12.py}
does not import that producer. It reconstructs allowable
products from associativity with the six band elements,
assigns the undetermined cells in reverse fixed order,
and never propagates an unknown outside product. It
checks the entire literal set, with $258$ search nodes
and $300612$ evaluated associativity equations, and
finds exactly the same $27$ tables. It also independently
checks their identities, full associativity and exact
minimal ideals.

Every table has a rational row-projector certificate.
Put $A=\mathbb QM/\mathbb QK$ and let $V$ be its
two-dimensional row augmentation module. The certificate
specifies a literal pair $e,g$ satisfying
$e^2=e$, $g^2=e$, $eg=ge=g$, the idempotent
$\varepsilon=(e-g)/2$, a basis for the full left ideal
$A\varepsilon$, and a vector $v\in V$. The independent
checker verifies that the ideal has dimension two and
that $x\mapsto xv$ sends its basis to a basis of $V$.
It verifies this map's equivariance on every algebra
basis element. Thus $V\cong A\varepsilon$ is projective
in all $27$ cases. The $324$ explicit basis-equivariance
equations and all rational ranks pass.
Lemma~\ref{CP2:A-act-row-pd-one} now kills the products.

All tables, idempotents and bases are in
\path{converse/scope18/forced_three_row12.json}; its
digest and the independent completion receipt are in
\path{converse/scope18/forced_three_row12_checked.json}.
The calculation is newly generated from the finite
profiles and uses no inherited table or module certificate.
\end{proof}

\begin{lemma}[Proved: four outside elements and two columns]
\label{CP2:A-two-column-four-outside-bounded}
Let $M$ have a rectangular minimal ideal with two
columns and at most four outside elements. Every finite
right-act cone with bounded rational homology has zero
rational degree-two products, at any row or state count.
\end{lemma}
\begin{proof}
If its contracted algebra $A$ is local, use
Corollary~\ref{CP2:A-local-two-column-acts}.
Suppose $A$ is nonlocal. A point fixed by the whole
minimal band is fixed by $M$, since $zm=z(em)=z$
for a band idempotent $e$ and $em\in K$. Such a
point makes every second class relative to the monoid
vertex and gives zero products.
Without such a point, the projective orbit-kernel lemma
and constant-action compression give a monoid $N$
with the same homology as the act cone in every degree
at least three. In particular its homology is bounded.
Its contracted algebra is the same $A$, of dimension
at most four. The band--Tor formula and
Lemma~\ref{A:nonlocal-four} show that $H_4(BN;\mathbb Q)$
must vanish: otherwise it would recur in arbitrarily
high degrees. Thus the original cone also has zero
fourth rational homology, proving the claim.
\end{proof}

\begin{lemma}[Proved: a column quotient preserves act-cone $H^2$]
\label{CP2:A-column-quotient-act-H2}
Let $q:M\to M'$ collapse columns of a rectangular
minimal ideal, preserve its row set bijectively, and
leave every outside element distinct. Suppose the
right $M$-action on $S$ factors through this actual
monoid quotient. The induced functor between the
ordinary act cones induces an isomorphism on rational
second cohomology. No assertion about their higher
homology is needed here.
\end{lemma}
\begin{proof}
Write $R=\mathbb QM$, $R'=\mathbb QM'$ and
$L=\ker(\mathbb QS\to\mathbb Q)$. The projective
row modules $P=Re$ and $P'=R'q(e)$ identify under
$q$ with the same permutation module $\mathbb QI$;
their augmentation kernels identify with the same
module $V$. All of $P,V,L$ therefore factor through
$R'$, and their $R$-linear maps are precisely their
$R'$-linear maps. Dualizing the reduced act-cone
Tor description and using $0\to V\to P\to\mathbb Q\to0$
identifies second cohomology with
\[
 \operatorname{coker}\bigl(
 \operatorname{Hom}_R(P,DL)\longrightarrow
 \operatorname{Hom}_R(V,DL)\bigr),
 \qquad D=\operatorname{Hom}_{\mathbb Q}(-,\mathbb Q).
\]
The identical presentation over $R'$ proves the stated
isomorphism. Naturality of the augmentation sequence
identifies it with the pullback of the actual quotient
functor, so it preserves cup products.
\end{proof}

\begin{lemma}[Proved: three columns, three outside elements, five states]
\label{CP2:A-three-by-three-five-state}
A five-state ordinary right-act cone over a monoid
with minimal band $3\times3$ and at most three outside
elements cannot have type $\CP^2$.
\end{lemma}
\begin{proof}
A band-fixed state is globally fixed and gives zero
products. Otherwise every core orbit in $SK$ has
two or three states. Each such orbit is individually
$M$-invariant: if $x=y(i,j)$, then
$xm=y(i,\tau_m(j))$ lies in the same orbit.
A projective-plane cone has connected action category,
since its first rational homology is the reduced
zeroth homology of that category. Thus two core orbits
cannot exhaust $S$.

With one core orbit every band element acts constantly
on all states: it is constant on that orbit, and
$sk=(sk)k$ determines the same value for every $s$.
Compression then gives an actually equivalent monoid
of order at most $3+3\cdot5=18$, contradicting the
monoid minimum.

The only other possibility is two two-point core
orbits and one transient state $x$. The column maps
of the two orbits give surjections
$\pi_0,\pi_1:\{0,1,2\}\to\{0,1\}$, up to labelling
each orbit. If their two-part partitions differ,
the three identifications
$\pi_0(j)\sim\pi_1(j)$ join all four core states.
Indeed two distinct two-part partitions of a three-set
give a connected bipartite graph on their four parts.
The constant-action quotient therefore has just two
states, including $x$. Its degree-two products vanish,
and its $H^2$ surjection kills those on the original cone.

Suppose the partitions agree, and label both orbits
by a common quotient $\pi$ of the old column set.
Every right actor preserves each core orbit and obeys
$s_jm=s_{\tau_m(j)}$, so $\pi\tau_m$ depends only on
$\pi(j)$. Collapsing the band columns by $\pi$ is
therefore an actual congruence of $M$. Let $M'$ be
the quotient, with minimal band $3\times2$ and the
same at most three outside elements.

The action on $S$ factors through $M'$. This is clear
on the cores. For the remaining state $x$, the map
$j\mapsto x(i,j)$ lands, for each fixed row $i$, in
one core orbit and is its column map. Thus equal
$\pi$-values give the same result on $x$ as well.
Lemma~\ref{CP2:A-column-quotient-act-H2} makes the
quotient functor an $H^2$ isomorphism.
Every degree-two product of the $M'$ act cone vanishes
unconditionally by
Proposition~\ref{CP2:A-three-row-five-outside-acts}.
Pullback gives the contradiction. No boundedness
claim was transferred across this column quotient.
\end{proof}

\begin{theorem}[Proved: twelve corner elements and five states]
\label{CP2:A-twelve-five-state}
An ordinary act cone over a monoid of order at most
twelve with at most five states cannot have type
$\CP^2$.
\end{theorem}
\begin{proof}
Orders through eleven use
Theorem~\ref{CP2:A-bounded-eleven-acts}; at order twelve,
at most four states are already excluded. A possible
nonzero square requires at least three minimal rows
and two columns. A nontrivial maximal group would
cost at least thirteen elements, so the ideal is a band.
At order twelve the only shapes are $3\times2$ with
six outside elements, $4\times2$ with four,
$5\times2$ with two, and $3\times3$ with three.
The first is the fresh projective-row calculation,
the next two are the four-outside bounded lemma,
and the last is the preceding five-state lemma.
Every possibility is excluded.
\end{proof}

\begin{corollary}[Proved: ordinary act cones through eighteen arrows]
\label{CP2:A-ordinary-act-eighteen}
No ordinary finite-monoid act cone with at most eighteen
total arrows, identities counted, has type $\CP^2$.
\end{corollary}
\begin{proof}
Write its size as $m+s+1\leq18$. For $m\leq11$
use the bounded-eleven theorem. If $m=12$, then
$s\leq5$ and the preceding theorem applies. For
$m=13$, $s\leq4$ and
Theorem~\ref{CP2:A-four-state-fourteen} applies.
For $m\geq14$ there are at most three states, excluded
by the universal one- and two-state obstructions and
the exact three-state minimum of twenty-two arrows.
\end{proof}
